\documentclass[12pt,a4paper]{article}

\usepackage[utf8]{inputenc}
\usepackage[T1]{fontenc}
\usepackage[polish]{babel}
\usepackage{geometry}
\usepackage{booktabs}
\usepackage{longtable}
\usepackage{listings}
\usepackage{xcolor}
\usepackage{hyperref}
\usepackage{titlesec}
\usepackage{parskip}
\usepackage{graphicx}
\usepackage{microtype}
\usepackage{fancyhdr}
\usepackage{enumitem}

\geometry{margin=2.5cm}

% --- kolory do listingów ---
\definecolor{codebg}{HTML}{1E1E1E}
\definecolor{codegreen}{HTML}{6A9955}
\definecolor{codekw}{HTML}{569CD6}
\definecolor{codestr}{HTML}{CE9178}
\definecolor{codefg}{HTML}{D4D4D4}
\definecolor{accentblue}{HTML}{2E6DA4}

% --- styl listingów ---
\lstdefinestyle{ada}{
  backgroundcolor=\color{codebg},
  basicstyle=\ttfamily\footnotesize\color{codefg},
  keywordstyle=\color{codekw}\bfseries,
  commentstyle=\color{codegreen},
  stringstyle=\color{codestr},
  breaklines=true,
  frame=single,
  rulecolor=\color{gray!40},
  framesep=4pt,
  xleftmargin=8pt,
  xrightmargin=4pt,
  tabsize=4,
  showstringspaces=false,
  morekeywords={procedure,function,is,begin,end,if,then,else,elsif,while,loop,
                for,in,return,array,of,Put_Line},
}

\lstdefinestyle{c}{
  backgroundcolor=\color{codebg},
  basicstyle=\ttfamily\footnotesize\color{codefg},
  keywordstyle=\color{codekw}\bfseries,
  commentstyle=\color{codegreen},
  stringstyle=\color{codestr},
  breaklines=true,
  frame=single,
  rulecolor=\color{gray!40},
  framesep=4pt,
  xleftmargin=8pt,
  xrightmargin=4pt,
  tabsize=4,
  showstringspaces=false,
  language=C,
}

\lstdefinestyle{bnf}{
  backgroundcolor=\color{gray!8},
  basicstyle=\ttfamily\footnotesize,
  breaklines=true,
  frame=single,
  rulecolor=\color{gray!40},
  framesep=4pt,
  xleftmargin=8pt,
  tabsize=4,
}

% --- nagłówek/stopka ---
\pagestyle{fancy}
\fancyhf{}
\rhead{\textcolor{gray}{\small Konwerter Ada $\rightarrow$ C}}
\lhead{\textcolor{gray}{\small Dokumentacja techniczna}}
\cfoot{\thepage}
\renewcommand{\headrulewidth}{0.4pt}

% --- kolory linków ---
\hypersetup{
  colorlinks=true,
  linkcolor=accentblue,
  urlcolor=accentblue,
}

% ===================================================================
\begin{document}

\begin{titlepage}
  \centering
  \vspace*{2cm}
  {\Huge\bfseries Konwerter kodu Ada~$\rightarrow$~C\par}
  \vspace{0.5cm}
  {\Large Dokumentacja techniczna\par}
  \vspace{1.5cm}
  \rule{\linewidth}{0.4pt}
  \vspace{1cm}

  \begin{tabular}{ll}
    \textbf{Autorzy:}   & Natalia Poźniak \\
                        & Małgorzata Poskróbek \\[6pt]
    \textbf{Język impl.:} & Java 17 \\
    \textbf{Generator parsera:} & ANTLR4 4.13.1 \\
    \textbf{Framework web:} & Spring Boot 3.2.3 \\
    \textbf{System budowania:} & Maven \\
  \end{tabular}

  \vspace{1cm}
  \rule{\linewidth}{0.4pt}
  \vfill
  {\large AGH Akademia Górniczo-Hutnicza\par}
  {\large \the\year\par}
\end{titlepage}

\tableofcontents
\newpage

% ===================================================================
\section{Opis projektu}

Projekt stanowi \textbf{transpiler} (source-to-source compiler) tłumaczący kod źródłowy
napisany w języku \textbf{Ada} na ekwiwalentny kod w języku~\textbf{C}.
Udostępnia dwa tryby działania:

\begin{itemize}
  \item \textbf{Aplikacja webowa} -- interfejs przeglądarkowy oparty na Spring Boot,
        pozwalający wpisać kod Ady i natychmiast uzyskać wynik konwersji,
  \item \textbf{Tryb CLI} -- uruchamiany bezpośrednio z wiersza poleceń;
        przetwarza wskazane pliki \texttt{.adb} i zapisuje wyniki do katalogu \texttt{out/}.
\end{itemize}

% ===================================================================
\section{Technologie}

\begin{table}[h!]
\centering
\begin{tabular}{lll}
\toprule
\textbf{Warstwa}       & \textbf{Technologia}                   & \textbf{Wersja} \\
\midrule
Język implementacji    & Java                                   & 17              \\
Framework webowy       & Spring Boot (spring-boot-starter-web)  & 3.2.3           \\
Generator parsera      & ANTLR4                                 & 4.13.1          \\
System budowania       & Apache Maven                           & ---             \\
Frontend               & HTML + Bootstrap 5.3 + vanilla JS      & ---             \\
\bottomrule
\end{tabular}
\caption{Stos technologiczny projektu}
\end{table}

% ===================================================================
\section{Architektura systemu}

Pipeline przetwarzania kodu przebiega przez cztery etapy:

\begin{enumerate}
  \item \textbf{Leksykalna analiza} -- \texttt{AdaLexer} (generowany z \texttt{AdaLexer.g4})
        zamienia strumień znaków na sekwencję tokenów.
  \item \textbf{Analiza składniowa} -- \texttt{AdaParser} (generowany z \texttt{AdaParser.g4})
        buduje drzewo rozbioru (\emph{ParseTree}).
  \item \textbf{Analiza semantyczna i generowanie kodu} -- \texttt{AdaToCVisitor}
        przechodzi drzewo metodą \emph{Visitor} i produkuje kod C.
  \item \textbf{Wyjście} -- tekst kodu C zwracany do wywołującego
        (plik \texttt{.c} w trybie CLI, odpowiedź JSON w trybie web).
\end{enumerate}

\begin{lstlisting}[style=bnf, caption={Schemat pipeline'u}]
Kod Ada (tekst)
      |
      v
  AdaLexer          <- tokenizacja (AdaLexer.g4)
      |
      v
  AdaParser         <- analiza skladniowa (AdaParser.g4)
      |  ParseTree
      v
  AdaToCVisitor     <- semantyka + generowanie kodu C
      |
      v
  Kod C (tekst)
\end{lstlisting}

% ===================================================================
\section{Opis klas}

\subsection{\texttt{Application}}

Punkt startowy aplikacji webowej. Uruchamia kontener Spring Boot na porcie 8080.

\begin{lstlisting}[style=c, caption={Klasa Application}]
@SpringBootApplication
public class Application {
    public static void main(String[] args) {
        SpringApplication.run(Application.class, args);
    }
}
\end{lstlisting}

\subsection{\texttt{Main}}

Punkt startowy trybu CLI. Przyjmuje listę plików \texttt{.adb} jako argumenty,
dla każdego wywołuje \texttt{Compiler.compileFile()} i zapisuje wynik do
katalogu \texttt{out/}. Błędy syntaktyczne i semantyczne trafiają na
\texttt{stderr}; program kończy się kodem~\texttt{2} w~przypadku niepowodzenia.

\subsection{\texttt{Compiler}}

Fasada enkapsulująca pipeline ANTLR4 dla trybu plikowego.
Odpowiada za:
\begin{itemize}
  \item wczytanie pliku (\texttt{CharStreams.fromFileName}),
  \item konfigurację leksera i parsera z własnym \texttt{ErrorListener}
        rzucającym \texttt{ParseCancellationException} przy błędzie składni,
  \item uruchomienie \texttt{AdaToCVisitor},
  \item udostępnienie listy błędów semantycznych przez
        \texttt{getSemanticErrors()} i \texttt{hasSemanticErrors()}.
\end{itemize}

\subsection{\texttt{CompilerController}}

Kontroler REST Spring Boot (\texttt{@RestController}).

\begin{itemize}
  \item \textbf{Endpoint:} \texttt{POST /api/compile}
  \item \textbf{Wejście:} JSON \texttt{\{"code": "..."\}}
  \item \textbf{Wyjście (sukces):} \texttt{\{"success": true, "cCode": "..."\}}
  \item \textbf{Wyjście (błąd):} \texttt{\{"success": false, "errors": [...]\}}
\end{itemize}

Pipeline działa na łańcuchu znaków (\texttt{CharStreams.fromString}),
bez operacji wejścia/wyjścia na dysku.

\subsection{\texttt{AdaToCVisitor}}

Główna klasa logiki transpilatora -- rozszerza wygenerowaną klasę
\texttt{AdaParserBaseVisitor<String>}. Każda metoda \texttt{visitXxx()}
odpowiada regule gramatycznej i zwraca fragment kodu C jako \texttt{String}.

\subsubsection{Zarządzanie zakresami}

Klasa utrzymuje stos (\texttt{Deque}) dwóch równoległych struktur:
\begin{itemize}
  \item \texttt{scopeTypes} -- \texttt{Map<String, CType>} mapująca
        nazwę zmiennej na jej typ,
  \item \texttt{declaredVars} -- \texttt{Set<String>} zmiennych
        zadeklarowanych w bieżącym zakresie.
\end{itemize}
Metody \texttt{pushScope()} / \texttt{popScope()} są wywoływane na wejściu
i wyjściu z każdej procedury i funkcji.

\subsubsection{System typów}

Wewnętrzny enum \texttt{CType}:

\begin{lstlisting}[style=c, caption={Enum CType}]
private enum CType {
    INT("int"),
    DOUBLE("double"),
    BOOL("bool");
}
\end{lstlisting}

Mapowanie typów Ada $\rightarrow$ C:

\begin{table}[h!]
\centering
\begin{tabular}{ll}
\toprule
\textbf{Typ Ada}                  & \textbf{Typ C} \\
\midrule
\texttt{Integer}, \texttt{Natural}, \texttt{Positive} & \texttt{int}    \\
\texttt{Float}                    & \texttt{double} \\
\texttt{Boolean}                  & \texttt{bool}   \\
inny (nieznany)                   & błąd (\texttt{RuntimeException}) \\
\bottomrule
\end{tabular}
\caption{Mapowanie typów}
\end{table}

Metoda \texttt{inferType()} przechodzi rekurencyjnie przez węzły
\texttt{expression $\rightarrow$ term $\rightarrow$ factor} i wyznacza typ
wynikowy wyrażenia.

\subsubsection{Generowanie kodu}

Wcięcia zarządzane są zmienną \texttt{indentLevel} (4 spacje na poziom).
Metoda \texttt{indent()} generuje odpowiedni łańcuch białych znaków.

% ===================================================================
\section{Obsługiwane konstrukcje języka Ada}

\begin{longtable}{ll}
\toprule
\textbf{Konstrukcja Ada} & \textbf{Odpowiednik C} \\
\midrule
\endhead
\texttt{procedure P is \ldots\ end P;} & \texttt{void P() \{ \ldots\ \}} \\
\texttt{procedure Main is \ldots} & \texttt{int main() \{ \ldots\ return 0; \}} \\
\texttt{function F return T is \ldots} & \texttt{T F() \{ \ldots\ \}} \\
\texttt{X : Integer := E;} & \texttt{int X = E;} (deklaracja ze typem) \\
\texttt{X : Float;} & \texttt{double X;} (deklaracja bez wartości) \\
\texttt{X : array (A..B) of Integer;} & \texttt{int X[B+1];} (tablica 1-indeksowana) \\
\texttt{X := E;} & \texttt{X = E;} \\
\texttt{if \ldots\ then \ldots\ elsif \ldots\ else \ldots\ end if;} & \texttt{if (\ldots) \{ \ldots\ \} else if \{ \ldots\ \} else \{ \ldots\ \}} \\
\texttt{while \ldots\ loop \ldots\ end loop;} & \texttt{while (\ldots) \{ \ldots\ \}} \\
\texttt{for I in A .. B loop \ldots} & \texttt{for (int I = A; I <= B; I++) \{ \}} \\
\texttt{return E;} & \texttt{return E;} \\
\texttt{Put\_Line(E);} & \texttt{printf("\%d$\backslash$n", E)} / \texttt{\%f} \\
\texttt{A(I)}, \texttt{A(I,J)} & \texttt{A[I]}, \texttt{A[I][J]} \\
\texttt{=}, \texttt{/=} & \texttt{==}, \texttt{!=} \\
\bottomrule
\caption{Obsługiwane konstrukcje}
\end{longtable}

% ===================================================================
\section{Analiza semantyczna}

Klasa \texttt{AdaToCVisitor} przeprowadza analizę semantyczną
w~trakcie przechodzenia drzewa. Wykryte błędy są zbierane na liście
\texttt{semanticErrors} i~zwracane po zakończeniu kompilacji.
Każdy komunikat zawiera numer linii w~oryginalnym kodzie Ada
(np.\ \texttt{Linia 5: Undeclared variable used: x}).

\begin{table}[h!]
\centering
\begin{tabular}{ll}
\toprule
\textbf{Rodzaj błędu} & \textbf{Przykład wyzwalający} \\
\midrule
Użycie niezadeklarowanej zmiennej & \texttt{x := y + 1} (y nieznane) \\
Niezgodność typów w przypisaniu & przypisanie \texttt{double} do \texttt{int} \\
Operacja arytmetyczna na \texttt{bool} & \texttt{True + 1} \\
Dzielenie przez zero (statyczne) & \texttt{x := a / 0;} \\
Nieprawidłowy typ indeksu tablicy & \texttt{A(1.5)} \\
Indeks tablicy poza zakresem & \texttt{A(10)} przy rozmiarze 5 \\
Nieznany typ Ada & \texttt{x : Integre := 1;} \\
Pusta funkcja (brak \texttt{return}) & \texttt{function F return Integer is begin end F;} \\
\bottomrule
\end{tabular}
\caption{Wykrywane błędy semantyczne}
\end{table}

% ===================================================================
\section{Gramatyka}

Gramatyka jest podzielona na dwa pliki ANTLR4.

\subsection{Lekser -- \texttt{AdaLexer.g4}}

Lekser obsługuje słowa kluczowe \emph{case-insensitive} (Ada jest językiem
nierozróżniającym wielkości liter) za pomocą fragmentów literowych:

\begin{lstlisting}[style=bnf, caption={Fragment AdaLexer.g4}]
PROCEDURE : P R O C E D U R E ;
FUNCTION  : F U N C T I O N ;
...
FLOAT      : [0-9]+ '.' [0-9]+ ;
INTEGER    : [0-9]+ ;
IDENTIFIER : [a-zA-Z_] [a-zA-Z0-9_]* ;
COMMENT    : '--' ~[\r\n]* -> skip ;
WS         : [ \t\r\n]+ -> skip ;
fragment A:[aA]; fragment B:[bB]; ...
\end{lstlisting}

\subsection{Parser -- \texttt{AdaParser.g4}}

\begin{lstlisting}[style=bnf, caption={Gramatyka (BNF -- skrót)}]
program          ::= subprogram_decl
subprogram_decl  ::= procedure_decl | function_decl

procedure_decl   ::= "procedure" IDENTIFIER "is"
                       var_decl_section
                     "begin"
                       proc_statement_list?
                     "end" IDENTIFIER ";"

function_decl    ::= "function" IDENTIFIER "return" IDENTIFIER "is"
                       var_decl_section
                     "begin"
                       func_statement_list?
                     "end" IDENTIFIER ";"

var_decl_section ::= var_decl*
var_decl         ::= IDENTIFIER ":" IDENTIFIER [ ":=" expression ] ";"
                   | IDENTIFIER ":" "array" "(" expression ".." expression ")" "of" IDENTIFIER ";"

proc_statement   ::= assignment | if_statement | while_statement
                   | for_statement | put_line_statement
func_statement   ::= assignment | if_statement | while_statement
                   | for_statement | return_statement

put_line_statement ::= "Put_Line" "(" expression ")" ";"

condition        ::= expression relational_op expression
relational_op    ::= "=" | "/=" | "<" | ">" | "<=" | ">="
expression       ::= term | expression ("+" | "-") term
term             ::= factor | term ("*" | "/") factor
factor           ::= "-" factor | IDENTIFIER | array_access
                   | INTEGER | FLOAT | "(" expression ")"
\end{lstlisting}

% ===================================================================
\section{Interfejs webowy}

Frontend serwowany jest jako plik statyczny z~\texttt{src/main/resources/static/index.html}.

\begin{itemize}
  \item \textbf{Lewy panel:} pole tekstowe z kodem Ada (monospacefont, ciemne tło).
  \item \textbf{Przycisk} \emph{Konwertuj do C}: wywołuje \texttt{POST /api/compile}
        przez Fetch API.
  \item \textbf{Prawy panel:} wygenerowany kod C (kolor niebieski) lub komunikat
        o~błędzie.
  \item \textbf{Sekcja błędów semantycznych:} lista błędów wyświetlana na czerwono
        poniżej paneli.
\end{itemize}

% ===================================================================
\section{Uruchomienie}

\subsection{Aplikacja webowa}

\begin{lstlisting}[style=bnf]
mvn spring-boot:run
\end{lstlisting}

Aplikacja dostępna pod adresem \texttt{http://localhost:8080}.

\subsection{Tryb CLI}

\begin{lstlisting}[style=bnf]
mvn package -DskipTests
java -jar target/ada-to-c-1.0.jar examples/factorial.adb
\end{lstlisting}

Wynik zapisywany jest w katalogu \texttt{out/factorial.c}.
W przypadku błędów program wypisuje komunikaty na \texttt{stderr}
i kończy działanie z kodem~\texttt{2}.

% ===================================================================
\section{Przykłady}

\subsection{Silnia (Ada)}

\begin{lstlisting}[style=ada, caption={examples/factorial.adb}]
procedure Factorial is
    n    : Integer := 5;
    fact : Integer := 1;
begin
    for i in 1 .. n loop
        fact := fact * i;
    end loop;
end Factorial;
\end{lstlisting}

\subsection{Wygenerowany kod C}

\begin{lstlisting}[style=c, caption={out/factorial.c}]
#include <stdio.h>
#include <stdbool.h>
#include <stdlib.h>

void Factorial() {
    int n = 5;
    int fact = 1;
    for (int i = 1; i <= n; i++) {
        fact = fact * i;
    }
}

int main() {
    Factorial();
    return 0;
}
\end{lstlisting}

\subsection{Przykład z błędem semantycznym}

\begin{lstlisting}[style=ada, caption={Kod Ada z błędem -- dzielenie przez zero}]
procedure DivByZero is
    x : Integer := 0;
    y : Integer := 0;
begin
    y := 10 / x;
end DivByZero;
\end{lstlisting}

Transpiler wykrywa dzielenie przez zero statycznie i zwraca błąd zamiast generować kod~C:

\begin{lstlisting}[style=bnf, caption={Odpowiedź JSON przy błędzie semantycznym}]
{
  "success": false,
  "errors": [
    "Linia 5: Blad semantyczny: Wykryto dzielenie przez zero za pomoca operatora '/'"
  ]
}
\end{lstlisting}

% ===================================================================
\section{Tabela tokenów}

\begin{longtable}{lll}
\toprule
\textbf{Token}   & \textbf{Wzorzec / Definicja}        & \textbf{Opis} \\
\midrule
\endhead
\texttt{PROCEDURE}  & \texttt{"procedure"} (case-insensitive) & deklaracja procedury \\
\texttt{FUNCTION}   & \texttt{"function"}                     & deklaracja funkcji \\
\texttt{IS}         & \texttt{"is"}                           & poczatek sekcji deklaracji \\
\texttt{BEGIN}      & \texttt{"begin"}                        & poczatek bloku instrukcji \\
\texttt{END}        & \texttt{"end"}                          & koniec bloku \\
\texttt{IF}         & \texttt{"if"}                           & instrukcja warunkowa \\
\texttt{THEN}       & \texttt{"then"}                         & czesc warunku if \\
\texttt{ELSE}       & \texttt{"else"}                         & alternatywa if \\
\texttt{ELSIF}      & \texttt{"elsif"}                        & dodatkowy warunek if \\
\texttt{WHILE}      & \texttt{"while"}                        & petla while \\
\texttt{LOOP}       & \texttt{"loop"}                         & blok petli \\
\texttt{FOR}        & \texttt{"for"}                          & petla for \\
\texttt{IN}         & \texttt{"in"}                           & operator zakresu w for \\
\texttt{RETURN}     & \texttt{"return"}                       & zwrot wartosci \\
\texttt{ARRAY}      & \texttt{"array"}                        & deklaracja tablicy \\
\texttt{OF}         & \texttt{"of"}                           & typ elementow tablicy \\
\texttt{ASSIGN}     & \texttt{:=}                             & przypisanie wartosci \\
\texttt{COLON}      & \texttt{:}                              & separator nazwy i typu \\
\texttt{RANGE}      & \texttt{..}                             & operator zakresu \\
\texttt{PLUS}       & \texttt{+}                              & dodawanie \\
\texttt{MINUS}      & \texttt{-}                              & odejmowanie \\
\texttt{MUL}        & \texttt{*}                              & mnozenie \\
\texttt{DIV}        & \texttt{/}                              & dzielenie \\
\texttt{EQ}         & \texttt{=}                              & rownosc \\
\texttt{NEQ}        & \texttt{/=}                             & nierownosc \\
\texttt{LT}         & \texttt{<}                              & mniejsze \\
\texttt{GT}         & \texttt{>}                              & wieksze \\
\texttt{LE}         & \texttt{<=}                             & mniejsze lub rowne \\
\texttt{GE}         & \texttt{>=}                             & wieksze lub rowne \\
\texttt{SEMICOLON}  & \texttt{;}                              & koniec instrukcji \\
\texttt{COMMA}      & \texttt{,}                              & separator \\
\texttt{LPAREN}     & \texttt{(}                              & nawias otwierajacy \\
\texttt{RPAREN}     & \texttt{)}                              & nawias zamykajacy \\
\texttt{INTEGER}    & \texttt{[0-9]+}                         & liczba calkowita \\
\texttt{FLOAT}      & \texttt{[0-9]+.[0-9]+}                  & liczba zmiennoprzecinkowa \\
\texttt{IDENTIFIER} & \texttt{[a-zA-Z\_][a-zA-Z0-9\_]*}      & nazwa zmiennej/procedury \\
\texttt{COMMENT}    & \texttt{-{}- \textasciitilde[\textbackslash r\textbackslash n]*} & komentarz (pomijany) \\
\texttt{WS}         & \texttt{[ \textbackslash t\textbackslash r\textbackslash n]+}  & biale znaki (pomijane) \\
\bottomrule
\caption{Tabela tokenow leksera}
\end{longtable}

\end{document}
